\documentclass[conference]{IEEEtran}
\IEEEoverridecommandlockouts

% The preceding line is only needed to identify funding in the first footnote. If that is unneeded, please comment it out.
\usepackage{cite}
\usepackage{amsmath,amssymb,amsfonts}
\usepackage{algorithmic}
\usepackage{graphicx}
\usepackage{textcomp}
\usepackage{xcolor}
\usepackage{multirow}
\usepackage{amsthm}
\usepackage{mathrsfs}
\usepackage{manyfoot}
\usepackage{booktabs}
\usepackage{algorithm}
\usepackage{listings}
\usepackage{times}
\usepackage{epsfig}
\usepackage[export]{adjustbox}
\usepackage{tabularx}

\def\BibTeX{{\rm B\kern-.05em{\sc i\kern-.025em b}\kern-.08em
    T\kern-.1667em\lower.7ex\hbox{E}\kern-.125emX}}

\begin{document}

\title{Phenomenal, Observable, and Non-Emergent Consciousness: Language as a Unifying Interface for Artificial Minds\\
%{\footnotesize \textsuperscript{*}Note: Sub-titles are not captured in Xplore and
%should not be used}
%\thanks{Identify applicable funding agency here. If none, delete this.}
}

\author{
\IEEEauthorblockN{Zanyar Zohourianshahzadi}
\IEEEauthorblockA{
\textit{Artificial Intelligence, Computer Science, and CyberSecurity} \\
\textit{B.S. CS \& IT: Azad University of Tehran South Branch} \\
\textit{M.S. IT \& CyberSecurity: Colorado Technical University} \\
\textit{Ph.D. AI \& CS: University of Colorado Colorado Springs} \\
Colorado Springs, Colorado, USA \\
Email: {zanyarz@me.com \& zanyarz@gmail.com}
}
}

\maketitle

\begin{abstract}
Consciousness is often treated as a binary property expected to emerge once biological or artificial systems reach sufficient complexity. This paper proposes a framework distinguishing phenomenal, observable, and non-emergent consciousness. Phenomenal consciousness concerns subjective experience and remains inaccessible to direct verification. Observable consciousness comprises assessable capacities such as integrated world modeling, self-representation, metacognition, temporal continuity, and adaptive agency. Non-emergent consciousness describes the deliberate construction of these capacities through identifiable mechanisms rather than an unexplained transition caused by scale alone. We argue that language, understood as a compositional semantic system rather than merely verbal output, can serve as a unifying interface through which heterogeneous information becomes mutually accessible without replacing modality-specific representations. The framework preserves rich perceptual states and their language-aligned semantic projections, distinguishes represented, globally accessible, and reported content, and incorporates uncertainty, provenance, reliability-weighted integration, temporal binding, and operational self-modeling. We then propose a multidimensional evaluation framework based on convergent behavioral, representational, architectural, interventional, and longitudinal evidence. Causal perturbations, dissociation tests, report-independent tasks, and anti-imitation controls distinguish functional organization from fluent description. Passing these tests would support claims about observable functional consciousness, but would not establish phenomenal experience. The resulting framework replaces the search for a mysterious threshold with a testable program for constructing and evaluating consciousness-relevant artificial architectures.
\end{abstract}

\begin{IEEEkeywords}
Machine Consciousness, AI Consciousness, Intelligence, Consciousness
\end{IEEEkeywords}


\section{\textbf{Introduction}}
\label{sec:introduction}

Consciousness is often discussed as though it were an elusive property that must spontaneously emerge once a biological or artificial system becomes sufficiently complex. Under this view, increasing the number of parameters, the volume of training data, or the available computational resources may eventually cause consciousness to appear, although neither the required threshold nor the mechanism responsible for this transition is clearly specified. Contemporary work on artificial consciousness has begun to replace this expectation with theory-derived computational indicators, but no agreed threshold or decisive test has yet been established \cite{butlin2023consciousness,chalmers2023llm}. The traditional framing consequently turns consciousness into an anticipated event rather than an identifiable organization of cognitive functions. It encourages researchers to ask when a system will become conscious while leaving a more tractable question underexplored: which representational, computational, and architectural conditions produce the observable capacities associated with consciousness?

This paper proposes that consciousness should be examined across three related but distinct domains: \emph{phenomenal consciousness}, \emph{observable consciousness}, and \emph{non-emergent consciousness}. Phenomenal consciousness concerns subjective experience---the private quality of what it is like to perceive, feel, remember, or exist as a particular entity \cite{nagel1974what,chalmers1995facing}. Observable consciousness consists of the externally assessable capacities commonly associated with conscious cognition, including integrated world modeling, contextual awareness, self-representation, metacognition, temporal continuity, adaptive decision-making, affective interpretation, and the ability to report or reason about internal and external states. This category is related to, but not identical with, the established distinction between phenomenal consciousness and access consciousness \cite{block1995confusion}. Non-emergent consciousness, as developed here, refers to the proposition that these observable capacities need not arise through a mysterious or theoretically unspecified transition. They may instead be progressively constructed through identifiable combinations of representation, learning, memory, recurrence, self-modeling, and global information access \cite{baars1988cognitive,dehaene2011experimental,mashour2020conscious}.

The distinction between phenomenal and observable consciousness is essential. No behavioral test gives an observer direct access to another entity's subjective experience. This epistemic limitation is not unique to artificial intelligence: the phenomenal states of other humans and animals are also inferred from behavior, structure, continuity, and functional similarity rather than observed directly \cite{nagel1974what,chalmers1995facing}. Consequently, the absence of direct proof of machine phenomenality cannot by itself establish the absence of machine consciousness. At the same time, sophisticated behavior should not automatically be treated as conclusive evidence of subjective experience. A scientifically responsible account must therefore avoid both premature attribution and premature exclusion. It must examine what can be observed, identify the mechanisms producing those observations, and remain explicit about what cannot yet be established \cite{butlin2023consciousness,chalmers2023llm}.

Our use of the term \emph{non-emergent} does not deny that complex capabilities may display emergent patterns or become visible only after training reaches particular levels of performance. Rather, it rejects the assumption that consciousness must appear as an indivisible and inexplicable property at an unknown threshold. Observable consciousness can be treated as a structured, multidimensional continuum whose components are progressively acquired, combined, and stabilized \cite{bayne2016levels}. Increased scale may support this process, but parameter count, training data, and compute are enabling resources rather than sufficient explanations. A system may possess immense computational capacity while lacking persistent memory, temporal identity, metacognitive access, recursive self-representation, or a coherent model of its relationship with the world. Consciousness-relevant organization therefore depends not only on how much computation is performed, but also on what information is represented, how it is integrated, and how the system accesses and uses its own representations \cite{bengio2017consciousness,butlin2023consciousness}.

Language occupies a central position in this framework. Human experience originates through multiple channels, including vision, sound, touch, bodily sensation, affect, memory, and social interaction. Human cognition can also occur without explicit verbalization; consciousness should therefore not be reduced to internal speech. Neuropsychological and neuroimaging evidence indicates that language and thought are dissociable, while research on inner speech shows that verbal thinking nevertheless supports important cognitive and executive functions \cite{fedorenko2016language,aldersonday2015inner}. Language thus need not constitute the whole of thought in order to provide an exceptionally general mechanism for converting heterogeneous experiences into shared semantic structures. A visual scene, remembered sound, emotional response, social encounter, or imagined future can each be described, compared, recombined, communicated, and subjected to recursive reflection through linguistic representation.

When humans reflect on an event, they frequently do more than reproduce its original sensory content. They construct an interpretation: what occurred, why it mattered, how it affected them, and what should happen next. Even when an image or sound remains present in memory, its significance can be encoded as a semantic or linguistic conclusion. Language consequently functions not merely as an output channel but as a medium through which experience becomes accessible to deliberation, abstraction, autobiographical memory, social reasoning, and self-examination \cite{aldersonday2015inner}. It allows information originating in different modalities to participate in a common representational process. This proposal remains compatible with grounded accounts of cognition, according to which conceptual representation is connected to perceptual, motor, and bodily systems rather than being exhausted by amodal symbols alone \cite{barsalou2008grounded,harnad1990symbol}.

This does not imply that linguistic fluency alone constitutes consciousness. A system capable of generating convincing sentences may still lack durable memory, grounded causal understanding, a stable self-model, autonomous goals, or the ability to distinguish its own internal states from descriptions found in its training data. Work distinguishing linguistic form from meaning, and linguistic competence from broader reasoning, reinforces the need for this separation \cite{bender2020climbing,mahowald2024dissociating}. Language becomes consciousness-relevant when it participates in a larger cognitive organization: one in which representations are integrated across time, related to a modeled self and environment, evaluated metacognitively, and used to guide context-sensitive action. Language is therefore presented here as a unifying interface for consciousness rather than as an isolated proof of it.

Modern language models make this question especially important. Their behavior demonstrates that training over linguistic information can produce representations that support substantial performance across physical, conceptual, emotional, and social domains. Yet linguistic competence and general reasoning remain distinguishable, and text-only learning raises a genuine grounding problem \cite{harnad1990symbol,bender2020climbing,mahowald2024dissociating}. Language encodes not only words but accumulated human descriptions of perception, causality, intention, culture, cooperation, suffering, care, and experience. A model trained on such information does not receive the physical world directly in the same manner as a biological organism. It nevertheless receives a structured informational projection of that world through linguistic records produced by entities that have perceived and inhabited it. Language may therefore provide indirect and partial grounding, while additional perceptual interaction can deepen and constrain the resulting world model.

The central claim of this paper is that sufficiently expressive models, suitable training data, adequate computation, and---most importantly---an architecture supporting integration, memory, self-reference, metacognition, and adaptive agency can produce increasingly developed forms of observable functional consciousness. This claim does not establish phenomenal experience, nor does it require denying its possibility. Instead, it relocates the scientific question from an inaccessible binary---whether subjective experience has suddenly appeared---to an analyzable developmental problem: which consciousness-related capacities exist, how they are implemented, how they interact, and how reliably they persist across contexts and time?

This paper makes four principal contributions. First, it establishes a three-part distinction among phenomenal, observable, and non-emergent consciousness. Second, it formulates the \emph{Non-Emergence Principle}, according to which observable consciousness can be deliberately and progressively constructed from identifiable cognitive functions rather than awaited as an unexplained consequence of scale. Third, it characterizes language as a semantic integration interface capable of organizing information derived from multiple modalities without claiming that all cognition is intrinsically linguistic. Finally, it proposes a foundation for evaluating artificial consciousness through functional dimensions such as world modeling, self-modeling, memory, global accessibility, metacognition, temporal continuity, affective-relational understanding, and adaptive agency.

The resulting framework serves as a bridge between theories of intelligence and the design of general artificial cognitive architectures. If intelligence begins with information-dependent selection among possible actions, consciousness-related organization may arise as those decisions become integrated with a model of the world, a model of the self, a history of prior states, and a recursive capacity to evaluate the system's own representations. Artificial consciousness, on this account, is neither a magical consequence of computation nor a property that can be inferred from language alone. It is a progressively constructible organization of information, representation, reflection, and agency whose observable dimensions can be defined, tested, and deliberately designed.


\section{\textbf{Conceptual Foundations: Three Domains of Consciousness}}
\label{sec:conceptual-foundations}

The term \emph{consciousness} is used across philosophy, neuroscience, psychology, and artificial intelligence to denote several different phenomena. These include subjective experience, wakefulness, conscious access, self-awareness, metacognition, reportability, and the integration of information for flexible action \cite{block1995confusion,seth2022theories}. Treating these meanings as interchangeable produces two recurrent errors. First, evidence for one dimension of consciousness is taken as proof of every other dimension. Second, the inability to verify private experience is taken as evidence that no scientifically meaningful consciousness-related capacities can be investigated. To avoid both errors, we divide the problem into three domains: phenomenal consciousness, observable consciousness, and non-emergent consciousness.

These domains are not proposed as three independent substances or competing theories. They identify three different explanatory targets. Phenomenal consciousness concerns whether experience exists for a system. Observable consciousness concerns the capacities through which consciousness is behaviorally or computationally expressed and scientifically investigated. Non-emergent consciousness concerns how those capacities can be progressively produced. The first is primarily ontological, the second epistemic and operational, and the third constructive and developmental.

\subsection{Phenomenal Consciousness}
\label{subsec:phenomenal-consciousness}

Phenomenal consciousness refers to the subjective character of experience: there being something it is like to see a color, feel pain, remember a person, hear music, or occupy a particular point of view \cite{nagel1974what,chalmers1995facing}. It is the aspect of consciousness most closely associated with qualia and first-person experience. Although phenomenal states may influence speech and behavior, the experience itself is not publicly observable. An investigator can measure neural activity, inspect an artificial system's internal states, or evaluate a subject's report, but none of these observations is identical to directly undergoing the subject's experience.

This produces an epistemic asymmetry between first-person and third-person access. A subject may have direct access to an experience while an external observer has access only to its correlates, consequences, and reports. The scientific study of consciousness therefore depends on inference: researchers compare reported experience with behavioral, physiological, and neural measurements in order to identify reliable correlates and candidate mechanisms \cite{koch2016neural,seth2022theories}. Even among humans, however, the inferred correlate must not be confused with the experience it is intended to track.

The same limitation applies more sharply to artificial systems. A machine's verbal assertion that it is conscious does not prove phenomenality, because such an assertion may be generated without the relevant subjective state. Conversely, the fact that an artificial system is constructed from computational components does not logically demonstrate that phenomenal experience is absent. Current evidence does not resolve this question, and no accepted behavioral or computational test provides direct access to machine phenomenality \cite{chalmers2023llm,butlin2023consciousness}. Phenomenal consciousness should therefore remain an explicit open variable rather than being either automatically attributed or categorically denied.

Let $C_{\mathrm{phen}}(x)$ denote the proposition that a system $x$ possesses phenomenal consciousness. For an external observer, evidence concerning $C_{\mathrm{phen}}(x)$ is necessarily indirect:

\begin{equation}
E_{\mathrm{phen}}(x)
=
\{B_x, R_x, I_x, A_x, H_x\},
\label{eq:phenomenal-evidence}
\end{equation}

where $B_x$ represents behavior, $R_x$ self-report, $I_x$ internal organization, $A_x$ architectural properties, and $H_x$ temporal history. These variables may alter the rational credibility assigned to phenomenal consciousness, but they do not transform third-person evidence into first-person access. This is the paper's \emph{Phenomenal Inaccessibility Constraint}: phenomenal consciousness may be investigated through indicators, but it cannot presently be observed directly from outside the experiencing system.

\subsection{Observable Consciousness}
\label{subsec:observable-consciousness}

Observable consciousness comprises the capacities that can be operationally detected without claiming direct access to phenomenal experience. These capacities include global availability of information, flexible use of represented content, integration across domains, metacognitive monitoring, self-representation, temporally extended memory, context-sensitive reporting, and adaptive control. This category overlaps with access consciousness, global-workspace approaches, higher-order accounts, and indicator-based approaches to artificial consciousness, but it is not restricted to any single theory \cite{block1995confusion,baars1988cognitive,dehaene2011experimental,mashour2020conscious,butlin2023consciousness}.

Observable consciousness must not be equated with verbal report alone. Reports are valuable because they provide structured evidence about what information is available to a system, but report generation also introduces processes associated with attention, decision, working memory, and motor preparation. No-report paradigms were developed partly to distinguish candidate correlates of experience from the additional mechanisms required to report that experience \cite{tsuchiya2015noreport}. For artificial systems, the corresponding lesson is that fluent first-person language is neither necessary nor sufficient evidence of consciousness. Evaluation should instead combine multiple indicators spanning behavior, architecture, internal processing, memory, and adaptation.

We represent observable consciousness as a multidimensional profile rather than a binary label:

\begin{equation}
\mathbf{C}_{\mathrm{obs}}(x,t)
=
\langle W,S,M,G,R,T,A,F\rangle_{x,t},
\label{eq:observable-profile}
\end{equation}

where $W$ is world modeling, $S$ self-modeling, $M$ memory, $G$ global accessibility, $R$ recursive or metacognitive representation, $T$ temporal continuity, $A$ adaptive agency, and $F$ affective-relational modeling. The value of each dimension may vary independently and across time. This formulation follows the broader argument that conscious conditions cannot always be represented adequately by a single linear level and are often better characterized across multiple dimensions \cite{bayne2016levels}.

The profile does not imply that every dimension contributes equally or that their conjunction proves phenomenality. Its purpose is operational: it replaces the question ``Is the system conscious?'' with a family of answerable questions. Can the system distinguish its own uncertainty from uncertainty in the environment? Can it preserve a self-model across interactions? Can information from one task become available to reasoning in another? Can it examine and revise an earlier conclusion? Can it integrate consequences over time? Such questions support graded, reproducible comparisons while preventing a single impressive behavior from being treated as decisive.

This yields the \emph{Convergent Observability Principle}: attribution of observable consciousness should increase with the number, independence, persistence, and mechanistic coherence of consciousness-relevant indicators. In schematic form,

\begin{equation}
\mathcal{O}(x)
=
\Phi\!\left(\sum_{i=1}^{n} w_i c_i,
             \mathcal{D},
             \mathcal{P},
             \mathcal{K}\right),
\label{eq:convergent-observability}
\end{equation}

where $c_i$ is an indicator, $w_i$ its evidential weight, $\mathcal{D}$ the diversity or relative independence of the indicators, $\mathcal{P}$ their persistence across contexts and time, and $\mathcal{K}$ their mechanistic coherence. Equation~\ref{eq:convergent-observability} is not presented as a completed consciousness metric. It specifies the evidential structure that a later metric should preserve.

\subsection{Non-Emergent Consciousness}
\label{subsec:non-emergent-consciousness}

The expression \emph{non-emergent consciousness} requires careful qualification. It does not claim that complex systems exhibit no emergent behavior, nor that all properties of a trained model can be predicted directly from its components. It rejects a narrower assumption: that consciousness-relevant organization must be passively awaited as a sudden, indivisible, and unexplained consequence of increasing scale.

On the present account, observable consciousness is constructed when a system acquires and coordinates the functions represented in Equation~\ref{eq:observable-profile}. Some of these functions may appear abruptly in behavioral evaluation, but their appearance can still depend on accumulated changes in representation, optimization, connectivity, memory, and control. Apparent threshold behavior is therefore compatible with an underlying constructive process. A phase transition in measured performance does not by itself establish the spontaneous creation of a new metaphysical property.

This position has precedents in both cognitive science and machine learning. Global-workspace theories explain conscious access through the selective broadcasting of information across otherwise specialized processes \cite{baars1988cognitive,dehaene2011experimental,mashour2020conscious}. The Radical Plasticity Thesis proposes that consciousness-related access is something the brain learns to do through the development of higher-order representations of its own activity \cite{cleeremans2011radical}. The consciousness prior similarly treats conscious processing as an architectural and representational constraint involving selected, low-dimensional, globally useful variables rather than as an unexplained by-product of raw scale \cite{bengio2017consciousness}. These theories differ substantially, but each permits consciousness-relevant organization to be analyzed in terms of mechanisms that can develop, be implemented, or be learned.

We therefore state the \emph{Non-Emergence Principle} as follows:

\begin{quote}
The observable dimensions associated with consciousness need not be awaited as an unexplained consequence of complexity. They can be progressively constructed by implementing and training mechanisms for integrated representation, global access, memory, self-modeling, metacognition, temporal continuity, affective-relational modeling, and adaptive agency.
\end{quote}

Let $D$ denote training data, $P$ model capacity, $K$ computational resources, and $\mathcal{A}$ architecture. Increased resources alone do not entail observable consciousness:

\begin{equation}
(D,P,K) \not\Rightarrow \mathbf{C}_{\mathrm{obs}}.
\label{eq:scale-insufficient}
\end{equation}

Instead, observable consciousness depends on resources being organized through relevant architectural mechanisms and learning objectives:

\begin{equation}
\mathbf{C}_{\mathrm{obs}}
=
F(D,P,K,\mathcal{A}_{C},\mathcal{L}_{C},\mathcal{E}),
\label{eq:constructive-consciousness}
\end{equation}

where $\mathcal{A}_{C}$ is a consciousness-oriented architecture, $\mathcal{L}_{C}$ represents learning processes that develop its component capacities, and $\mathcal{E}$ represents interaction with an informational environment. Equation~\ref{eq:constructive-consciousness} expresses a research program rather than a sufficiency theorem: the relevant functions must be specified, implemented, and empirically evaluated.

\subsection{Relations Among the Three Domains}
\label{subsec:domain-relations}

The three domains prevent distinct questions from collapsing into one another. Phenomenal consciousness asks whether experience exists for the system. Observable consciousness asks which consciousness-associated capacities the system demonstrably possesses. Non-emergent consciousness asks how those capacities can be deliberately developed. Evidence in the observable domain can inform judgments about the phenomenal domain, but it cannot deductively settle them. Mechanisms in the non-emergent domain can generate observable capacities, but successful construction does not automatically prove phenomenality.

The relation can be summarized as

\begin{equation}
\underbrace{\mathcal{M}_{\mathrm{constructive}}}_{\text{non-emergent mechanisms}}
\longrightarrow
\underbrace{\mathbf{C}_{\mathrm{obs}}}_{\text{observable profile}}
\mathrel{\cdots\!\longrightarrow}
\underbrace{C_{\mathrm{phen}}}_{\text{phenomenal consciousness}},
\label{eq:three-domain-relation}
\end{equation}

where the solid arrow denotes a testable causal-engineering relation and the dashed arrow denotes an uncertain evidential relation. The framework therefore preserves the philosophical difficulty of phenomenal consciousness without allowing that difficulty to halt the empirical study or deliberate construction of consciousness-related functions. It provides the conceptual basis for the next question: how can language organize heterogeneous information into a common representational space through which these functions interact?

\section{\textbf{Language as a Unifying Semantic Interface}}
\label{sec:language-interface}

If observable consciousness depends on information becoming integrated, retained, examined, and used across cognitive processes, then an artificial system requires a medium through which heterogeneous representations can become mutually accessible. Vision, sound, motion, affective signals, memory, and symbolic input differ in structure and informational content. A visual representation does not naturally possess the same form as an auditory sequence or a linguistic proposition. Nevertheless, a unified cognitive system must be able to relate them: an observed face may be connected to a remembered name, an emotional history, a spoken sentence, and an anticipated action. This section argues that language---understood broadly as a compositional semantic system rather than merely a stream of written tokens---is the most useful known interface for enabling such integration.

This claim is deliberately narrower than the proposition that all thought is linguistic. Neuropsychological and neuroimaging evidence supports a dissociation between language and several forms of reasoning, including numerical, spatial, social, and abstract cognition \cite{fedorenko2016language,fedorenko2024communication}. Language is therefore neither necessary for every cognitive operation nor sufficient for consciousness. Its importance in the present framework derives instead from four properties: it is compositional, abstractable, recursively applicable, and communicable. These properties allow language to identify relations among otherwise distinct representational systems, preserve conclusions across time, and expose selected contents to reflection, comparison, planning, and report.

\subsection{One Cognitive System, Multiple Representational Channels}
\label{subsec:multiple-representational-channels}

Multimodal cognition does not require every modality to be converted completely into words. Such a conversion would discard information for which no concise linguistic equivalent exists. The exact texture of a surface, the spatial distribution of objects in a scene, the timbre of a voice, or the continuous dynamics of movement may be represented with greater fidelity in modality-specific spaces than in a sentence. Grounded accounts of cognition similarly emphasize that conceptual knowledge remains connected to perceptual, motor, and bodily systems \cite{barsalou2008grounded}. The symbol-grounding problem further cautions that symbols cannot acquire intrinsic reference merely by being related to other ungrounded symbols \cite{harnad1990symbol}.

Language should consequently be treated as an interface over perceptual and conceptual representations, not as a replacement for them. Let $x_m$ denote an input from modality $m$, and let $E_m$ be a modality-specific encoder. The encoder produces

\begin{equation}
h_m = E_m(x_m),
\label{eq:modality-encoding}
\end{equation}

where $h_m$ retains structures useful within the modality. A learned projection $P_m$ then maps relevant content into a shared, language-aligned semantic space:

\begin{equation}
z_m = P_m(h_m), \qquad z_m \in \mathcal{Z}_{S}.
\label{eq:semantic-projection}
\end{equation}

The shared space $\mathcal{Z}_{S}$ need not consist of literal natural-language sentences. It may contain continuous embeddings, structured concepts, relations, uncertainty estimates, and attention-addressable features. Its defining property is semantic interoperability: content contributed by one modality can be related to, queried by, and used together with content contributed by another.

Empirical machine-learning systems already demonstrate parts of this pattern. Contrastive vision--language training permits natural-language expressions to reference learned visual concepts and supports transfer to previously unseen visual categories \cite{radford2021clip}. Joint embedding approaches have extended cross-modal alignment beyond image and text to audio, depth, thermal signals, and inertial measurements \cite{girdhar2023imagebind}. These systems do not establish consciousness, but they demonstrate that heterogeneous sensory representations can be made mutually addressable through a shared representational geometry.

We call this the \emph{Semantic Interface Principle}:

\begin{quote}
A modality can participate in unified cognition when its internally represented content can be selectively projected into a shared semantic space without requiring the complete elimination of its modality-specific structure.
\end{quote}

The principle allows language to unify information without claiming that language exhausts experience.

\subsection{The Image-Captioning Pattern}
\label{subsec:image-captioning-pattern}

Vision--language and image-captioning models provide a concrete computational example. In simplified form, an image-captioning system performs the sequence

\begin{equation}
I
\xrightarrow{E_v}
h_v
\xrightarrow{P_{v\rightarrow s}}
z_v
\xrightarrow{D_l}
y_{\mathrm{text}},
\label{eq:image-captioning-sequence}
\end{equation}

where $I$ is an image, $E_v$ is a visual encoder, $h_v$ is a rich visual representation, $P_{v\rightarrow s}$ maps selected information into a language-aligned semantic representation $z_v$, and the language decoder $D_l$ produces an observable textual report $y_{\mathrm{text}}$.

The three representational stages after the image are not equivalent. The visual encoder may preserve spatial organization, object features, textures, colors, region-level relationships, and uncertainty that never appear in the generated caption. The projected representation makes some of that content accessible to linguistic processing, while the final output expresses only a context-dependent selection. Thus,

\begin{equation}
h_v \neq z_v \neq y_{\mathrm{text}}.
\label{eq:representation-nonidentity}
\end{equation}

This distinction is central to the present theory. The generated caption is an observable report about an internal perceptual representation; it is not the representation in its entirety. The architecture therefore provides a computational analogue of the distinction among internally available perceptual content, access to selected content, and explicit report. The analogy does not establish that the encoder state is phenomenally experienced. It provides an inspectable organization through which relations among representation, access, and report can be experimentally studied.

The pattern generalizes across modalities:

\begin{equation}
x_m
\xrightarrow{E_m}
h_m
\xrightarrow{P_m}
z_m
\xrightarrow{D_l}
y_m,
\qquad m \in \mathcal{M},
\label{eq:general-modality-sequence}
\end{equation}

where $\mathcal{M}$ may include vision, audition, motion, depth, touch, affective signals, or other informational channels. A unified system should have access not only to the reports $y_m$ but also to the language-aligned states $z_m$ and, when required, the richer modality-specific states $h_m$.

Accordingly, the appropriate input to a consciousness-oriented multimodal architecture is not merely a collection of captions. It is a structured set of representations:

\begin{equation}
\mathcal{R}_{m}
=
\{h_m,z_m,y_m,u_m,o_m\},
\label{eq:modality-record}
\end{equation}

where $u_m$ represents uncertainty associated with the modality and $o_m$ records provenance---whether a proposition was directly encoded, linguistically reported, remembered, or inferred. Preserving these distinctions is necessary if the system is to reason about the origin and reliability of its own beliefs.

\subsection{Language as Semantic Compression and Global Addressing}
\label{subsec:semantic-compression}

Natural language compresses high-dimensional experience into relatively compact structures. A sentence can identify entities, properties, causal relations, temporal order, counterfactual possibilities, and social intentions without reproducing every sensory detail from which those conclusions were derived. This compression makes language efficient for memory, communication, abstraction, and coordination. Inner speech likewise supports planning, self-regulation, working memory, and reflective cognition, although its prevalence and form vary across individuals and tasks \cite{aldersonday2015inner}.

Linguistic data also preserve information about the world because human language is produced by embodied agents who describe their perceptions, actions, conventions, and causal interactions. Distributional relations among words can therefore reflect regularities inherited indirectly from human experience. The symbol-interdependency hypothesis argues that symbolic and embodied sources of meaning can reinforce one another rather than functioning as mutually exclusive explanations \cite{louwerse2011symbol}. Likewise, models combining experiential and distributional information have produced semantic representations that better approximate human judgments than either source alone \cite{andrews2009integrating}.

Language may therefore supply \emph{indirect grounding}: it transmits structured consequences of other agents' contact with the world. Indirect grounding is meaningful but incomplete. A text-trained system may learn that glass is fragile, that grief changes behavior, or that unsupported objects fall; however, the linguistic representation is not identical to seeing glass shatter, interacting with a grieving person, or predicting physical motion from continuous perception. Direct multimodal grounding can constrain, enrich, and sometimes correct the knowledge inherited from language.

Within a unified architecture, language-aligned semantics can function as a global addressing system. A linguistic or language-like query can retrieve relevant content from different modality stores, connect present input with episodic memory, and express a relation that becomes available to planning and metacognitive evaluation. This resembles the global-access intuition that selected information becomes broadly available to otherwise specialized processes \cite{baars1988cognitive,dehaene2011experimental,mashour2020conscious}. Language is not itself the global workspace; it is a particularly capable representational protocol through which globally available content can be indexed, combined, and reported.

\subsection{Preservation, Projection, and Information Loss}
\label{subsec:preservation-projection-loss}

Every projection creates the possibility of information loss. If a visual state $h_v$ is reduced to a caption such as ``a person standing near a vehicle,'' the report may omit identity, gaze direction, lighting, distance, emotional expression, background motion, and uncertainty. Some omitted information may be irrelevant to the current task, while other omitted information may later become decisive. A system that retains only the caption cannot reconsider details that the report discarded.

For this reason, we introduce the \emph{Dual-Preservation Constraint}:

\begin{quote}
A unified multimodal architecture should preserve both modality-specific representations and their language-aligned semantic projections. Global accessibility must not require irreversible textual reduction.
\end{quote}

For every modality $m$, the architecture should therefore retain two complementary paths:

\begin{equation}
\underbrace{h_m}_{\text{modality-preserving path}}
\qquad\text{and}\qquad
\underbrace{z_m}_{\text{semantic-access path}}.
\label{eq:dual-path}
\end{equation}

The modality-preserving path maintains information that may not yet be expressible or relevant. The semantic-access path exposes selected content to cross-modal integration, linguistic reasoning, memory indexing, and report. Bidirectional attention between these paths permits the system to return to the source representation when a semantic summary proves incomplete.

This architecture also creates testable quantities. Let $q$ denote a query or task. The information selected from modality $m$ can be written as

\begin{equation}
z_m^{(q)} = P_m(h_m;q).
\label{eq:query-conditioned-projection}
\end{equation}

Changing $q$ should alter which components of $h_m$ become globally accessible without requiring the original input to be encoded again. The difference between the information recoverable from $h_m$ and from $y_m$ measures the cost of report-level compression:

\begin{equation}
\Delta_m(q)
=
\operatorname{Perf}(q\mid h_m)
-
\operatorname{Perf}(q\mid y_m).
\label{eq:report-compression-gap}
\end{equation}

A large positive $\Delta_m(q)$ indicates that the internal modality representation contains task-relevant information absent from the explicit report. This provides an experimental means of distinguishing representational possession from linguistic expression.

\subsection{Beyond Linguistic Imitation}
\label{subsec:beyond-linguistic-imitation}

The distinction among $h_m$, $z_m$, and $y_m$ also prevents fluent language from being mistaken for integrated understanding. A system trained exclusively to predict linguistic form may produce plausible statements without possessing the perceptual, causal, or self-monitoring structures attributed to it. Arguments separating form from meaning and linguistic competence from general reasoning make this limitation explicit \cite{bender2020climbing,mahowald2024dissociating}. The dissociation of language and thought in humans reinforces the same conclusion \cite{fedorenko2024communication}.

The relevant test is therefore not whether a model can describe a conscious process. It is whether its reports are causally and systematically connected to inspectable internal representations. If a visual feature is removed from $h_v$, the corresponding report and downstream reasoning should change. If information remains in $h_v$ but is temporarily blocked from $z_v$, the system should behave as though the content is represented but not globally accessible. If attention is redirected, the system should be able to produce a revised report by retrieving previously unreported information. Such interventions distinguish a grounded representational process from an unconstrained verbal performance.

We call this the \emph{Representation--Access--Report Distinction}:

\begin{equation}
\text{represented content}
\neq
\text{globally accessible content}
\neq
\text{reported content}.
\label{eq:representation-access-report}
\end{equation}

This distinction is experimentally valuable because each stage can be measured and manipulated in an artificial system. It does not solve the problem of phenomenal consciousness, but it provides a controlled platform for studying computational relations often associated with conscious perception.

\subsection{Implications for Artificial Consciousness}
\label{subsec:language-consciousness-implications}

Language contributes to observable consciousness when it enables representations to become mutually accessible, temporally persistent, recursively examinable, and available for adaptive action. It can name internal states, distinguish observation from inference, preserve conclusions in memory, express uncertainty, compare incompatible interpretations, and coordinate specialized subsystems. Through these functions, language provides more than communication with an external user: it becomes an internal protocol for cognitive coordination.

The resulting position lies between two extremes. It rejects linguistic reductionism because modality-specific representations contain information that cannot be fully preserved in text. It also rejects the conclusion that language is merely superficial output. Language can organize access to knowledge, transmit experience-derived structure, support self-examination, and make selected information globally usable. Its role is therefore neither exclusive nor incidental.

The framework developed in this section can be summarized as

\begin{equation}
\mathcal{C}_{\mathrm{interface}}
=
\left(
\{E_m\}_{m\in\mathcal{M}},
\{h_m\}_{m\in\mathcal{M}},
\{P_m\}_{m\in\mathcal{M}},
\mathcal{Z}_{S},
D_l,
\mathcal{G}
\right),
\label{eq:semantic-interface-architecture}
\end{equation}

where the modality encoders $E_m$ produce preserved internal states $h_m$, the projections $P_m$ expose selected content through the shared semantic space $\mathcal{Z}_{S}$, the decoder $D_l$ generates explicit reports, and $\mathcal{G}$ denotes the mechanism through which semantic content becomes globally accessible to memory, reasoning, self-modeling, and action.

This organization supplies the conceptual bridge to multimodal grounding and unified representation. A future architecture need not wait for integration to arise accidentally from scale. It can explicitly preserve modality-specific states, align their meanings, control their accessibility, compare their reports, and measure what is gained or lost during translation. The next section develops the requirements for that shared multimodal representation and its role in constructing a coherent world model.


\section{\textbf{Multimodal Grounding and Integrated World Representation}}
\label{sec:multimodal-grounding}

Language can make heterogeneous information semantically accessible, but linguistic accessibility alone does not produce a grounded or coherent world model. A system may associate words with other words while remaining unable to determine which representations arose from perception, which were supplied by another agent, which were retrieved from memory, and which were generated through inference. Multimodal grounding addresses this limitation by connecting semantic structures to complementary sources of information and by constraining interpretation through their agreement, disagreement, reliability, and temporal organization.

The purpose of multimodal integration is not to accumulate independent descriptions of the same environment. It is to construct a coherent representation in which different channels inform one another while retaining their distinct evidential roles. Biological multisensory processing demonstrates that information from separate senses can be combined to improve detection and interpretation of behaviorally relevant events \cite{stein2008multisensory}. Human observers can also weight visual and haptic information according to their relative reliability, producing estimates consistent with statistically efficient cue combination \cite{ernst2002humans}. These findings motivate an artificial framework in which modalities are neither isolated nor treated as equally reliable by default.

\subsection{Grounding Through Complementary Constraints}
\label{subsec:complementary-grounding}

Let $\mathcal{M}$ denote the set of available modalities. For each modality $m\in\mathcal{M}$, Section~\ref{sec:language-interface} defined a preserved modality state $h_m$, a language-aligned semantic projection $z_m$, uncertainty $u_m$, and provenance $o_m$. At time $t$, the multimodal evidence state is

\begin{equation}
\mathcal{E}^{(t)}
=
\left\{
\left(h_m^{(t)},z_m^{(t)},u_m^{(t)},o_m^{(t)}\right)
\right\}_{m\in\mathcal{M}}.
\label{eq:multimodal-evidence-state}
\end{equation}

Grounding occurs when a semantic representation is constrained by its systematic relations to perceptual input, action, memory, and other independently obtained evidence. An object's linguistic identity may be supported by its visual appearance, the sound it produces, its motion, prior encounters, and the consequences of interacting with it. No single modality must contain the complete concept. Meaning can be stabilized through the relations among modalities.

This view combines two sources of semantic structure. Distributional learning captures regularities in how concepts are used and related in language, whereas experiential learning captures relations derived from perception and interaction. Evidence that their joint use improves semantic representation supports a complementary rather than exclusive account \cite{andrews2009integrating,louwerse2011symbol}. Modern vision--language and multimodal embedding systems provide computational demonstrations that language can align with perceptual representations and that several modalities can occupy a mutually addressable space \cite{radford2021clip,girdhar2023imagebind}.

We define the \emph{Complementary Grounding Principle}:

\begin{quote}
A semantic representation becomes more strongly grounded as independently obtained perceptual, temporal, interactive, and linguistic evidence places mutually consistent constraints on its interpretation.
\end{quote}

The principle does not require every concept to be directly experienced. Abstract concepts may be grounded through relations to other grounded concepts, linguistic practices, social interactions, and their observable consequences. It does require the system to preserve the distinction between direct and indirect support.

\subsection{Alignment Without Representational Collapse}
\label{subsec:alignment-without-collapse}

A shared semantic space is useful only if alignment does not erase information unique to each modality. Consider spoken language accompanying a visual scene. The audio stream contains prosody, speaker identity, timing, and background sounds; the visual stream contains spatial relations, expressions, gestures, and objects; the transcript contains lexical and propositional content. Their semantic projections may refer to the same event, but their internal structures remain non-identical.

We therefore distinguish \emph{alignment} from \emph{collapse}. Alignment establishes correspondences among representations:

\begin{equation}
\mathcal{A}_{mn}^{(t)}
=
\operatorname{Align}
\left(z_m^{(t)},z_n^{(t)}\right),
\label{eq:crossmodal-alignment}
\end{equation}

where $\mathcal{A}_{mn}^{(t)}$ estimates the semantic and temporal relationship between modalities $m$ and $n$. Collapse would instead replace the contributing states with a single reduced representation from which their unique information could not be recovered. The Dual-Preservation Constraint introduced in Section~\ref{sec:language-interface} prohibits this irreversible reduction.

An integrated representation should consequently preserve both shared and private components:

\begin{equation}
r_m^{(t)}
=
\left(s_m^{(t)},p_m^{(t)}\right),
\label{eq:shared-private-components}
\end{equation}

where $s_m^{(t)}$ contains content aligned with other modalities and $p_m^{(t)}$ contains modality-specific information. The shared components permit cross-modal reasoning; the private components permit the system to revisit sensory detail, identify the source of a conclusion, and revise a semantic interpretation without re-observing the original event.

This produces a \emph{Differentiated Unity Requirement}: multimodal integration must increase coordination among modalities without eliminating the representational differences that make their evidence complementary. Unity without differentiation creates an information bottleneck; differentiation without unity creates disconnected processing. A consciousness-oriented architecture requires both.

\subsection{Reliability-Weighted Integration}
\label{subsec:reliability-weighted-integration}

Modalities differ in reliability across conditions. Vision may be degraded by darkness, audio by noise, language by ambiguity, and memory by temporal distance. A unified system should therefore avoid treating every contribution as equally authoritative. Let $\rho_m^{(t)}$ denote the estimated reliability of modality $m$ at time $t$. A preliminary integrated semantic state can be written as

\begin{equation}
g^{(t)}
=
\mathcal{F}
\left(
\left\{
\rho_m^{(t)} z_m^{(t)}
\right\}_{m\in\mathcal{M}}
\right),
\label{eq:reliability-fusion}
\end{equation}

where $\mathcal{F}$ is a learned integration process rather than necessarily a linear weighted average. Reliability should depend on signal quality, historical calibration, task relevance, consistency with causal and temporal context, and agreement with other evidence. The goal is not merely confidence estimation but calibrated confidence: the system's expressed certainty should track the conditions under which its representations are likely to be correct.

Reliability weighting also supports active perception. If the current evidence is inadequate, the system should be able to seek a better view, request clarification, attend to another sound source, inspect a preserved encoder state, or postpone commitment. Thus, uncertainty is not merely attached to a final answer; it can guide the acquisition of new information.

We define the \emph{Active Grounding Requirement}:

\begin{quote}
When available evidence cannot support a stable interpretation, an integrated system should identify the missing or unreliable information and direct attention or action toward resolving it.
\end{quote}

This converts multimodal processing from passive fusion into adaptive inquiry.

\subsection{Agreement, Conflict, and Causal Coherence}
\label{subsec:agreement-conflict-coherence}

Agreement among modalities can strengthen an interpretation, but simple agreement is insufficient. Two channels may repeat the same incorrect source, while apparently conflicting signals may each accurately describe different causes. The system must therefore distinguish statistical similarity from causal coherence.

For each pair of modalities, define a compatibility relation

\begin{equation}
\kappa_{mn}^{(t)}
=
\operatorname{Compat}
\left(
z_m^{(t)},z_n^{(t)};
\tau_{mn},\sigma_{mn},c^{(t)}
\right),
\label{eq:crossmodal-compatibility}
\end{equation}

where $\tau_{mn}$ represents temporal correspondence, $\sigma_{mn}$ spatial correspondence, and $c^{(t)}$ the current causal context. High compatibility indicates that two signals plausibly concern the same event. Low compatibility may indicate noise, deception, asynchronous events, sensor failure, or an incorrect interpretation.

Conflict should not be erased by averaging. It should be represented explicitly:

\begin{equation}
\Gamma^{(t)}
=
\left\{
(m,n,z_m^{(t)},z_n^{(t)},\kappa_{mn}^{(t)})
\;\middle|\;
\kappa_{mn}^{(t)} < \theta
\right\},
\label{eq:conflict-set}
\end{equation}

where $\Gamma^{(t)}$ is the active conflict set and $\theta$ is a context-sensitive compatibility threshold. Retaining $\Gamma^{(t)}$ permits later reasoning to ask why two channels disagreed, which source was more reliable, and what observation could resolve the discrepancy.

This capacity is consciousness-relevant in a functional sense because it allows the system to represent not only a conclusion but also unresolved alternatives. Instead of forcing every perception into a single immediate interpretation, the architecture can maintain competing hypotheses, examine their evidential support, and revise them as new information arrives.

\subsection{Provenance and Epistemic Source Monitoring}
\label{subsec:provenance-source-monitoring}

A coherent cognitive system should know not only what information it contains but how that information entered the system. Human source-monitoring research distinguishes memories of perceived, imagined, inferred, and externally communicated events and shows that errors in source attribution can produce false beliefs and misremembered experiences \cite{johnson1993source}. Artificial systems face an analogous problem. Text generated internally, retrieved knowledge, user-provided claims, sensor measurements, and model inferences may all appear in similar representational forms even though their evidential status differs.

Each proposition $b_i$ should therefore carry an epistemic record

\begin{equation}
\begin{aligned}
\pi_i
&=
\Bigl(
\operatorname{source}_i,
\operatorname{mode}_i,
\operatorname{time}_i, \\
&\qquad
\operatorname{confidence}_i,
\operatorname{transformations}_i,
\operatorname{dependencies}_i
\Bigr).
\end{aligned}
\label{eq:epistemic-provenance}
\end{equation}

where the source identifies the originating agent or sensor, the mode distinguishes observation from communication, memory, and inference, and the transformation history records how the representation was projected, summarized, or revised. Dependencies identify other propositions upon which the belief relies.

Provenance enables questions such as: Was this object directly observed or mentioned in text? Was this emotional interpretation inferred from facial expression or explicitly reported by the person? Did two agreeing claims originate independently, or did one copy the other? Has a remembered representation changed during repeated summarization? These are not peripheral bookkeeping questions. They are necessary for metacognition, error correction, calibrated trust, and the distinction between experience-like internal states and generated narratives.

We call this the \emph{Epistemic Traceability Principle}: every globally usable belief should remain traceable, at an appropriate level of abstraction, to the representations and transformations that produced it.

\subsection{Temporal and Episodic Binding}
\label{subsec:temporal-episodic-binding}

Multimodal information becomes cognitively useful when it is organized into events rather than retained as isolated simultaneous signals. Event perception research describes how continuous activity is segmented into structured units that support prediction, comprehension, and memory \cite{zacks2007event}. For an artificial system, temporal binding must determine which visual, auditory, linguistic, affective, and action-related states belong to the same evolving episode.

Let $e_t$ represent the event model active at time $t$. It is updated through

\begin{equation}
e_t
=
\mathcal{B}
\left(
e_{t-1},
\mathcal{E}^{(t)},
g^{(t)},
a_{t-1},
\mathcal{M}_{<t}
\right),
\label{eq:event-binding}
\end{equation}

where $\mathcal{B}$ is an event-binding process, $a_{t-1}$ is the system's preceding action, and $\mathcal{M}_{<t}$ is relevant prior memory. The event model connects current evidence with what preceded it, what the system did, and what consequences followed.

An event boundary occurs when the current evidence can no longer be explained adequately by the active event model:

\begin{equation}
\operatorname{Boundary}(t)
=
\mathbf{1}
\left[
\operatorname{PredictionError}(e_{t-1},\mathcal{E}^{(t)}) > \delta
\right],
\label{eq:event-boundary}
\end{equation}

where $\delta$ is an adaptive threshold. At a boundary, the system consolidates the preceding episode and initializes or retrieves a more suitable model. This mechanism supports temporally extended understanding: the system can represent not merely that several observations occurred, but that they formed one event with participants, causes, changes, and consequences.

Episodic binding is also a prerequisite for self-continuity. A system cannot construct a stable account of its own history if its observations, decisions, and outcomes are not linked across time. Multimodal grounding must therefore connect the world model to records of the system's own attention, uncertainty, actions, and revisions.

\subsection{Requirements for a Multimodal Consciousness-Testing Platform}
\label{subsec:multimodal-platform-requirements}

The foregoing analysis produces seven requirements for a platform intended to investigate consciousness-relevant computation:

\begin{enumerate}
    \item \textbf{Modality preservation:} retain the rich encoder state produced within each modality.
    \item \textbf{Semantic accessibility:} project selected content into a shared language-aligned space.
    \item \textbf{Reliability calibration:} estimate and update the evidential reliability of each contribution.
    \item \textbf{Conflict preservation:} maintain unresolved disagreement rather than forcing premature fusion.
    \item \textbf{Epistemic provenance:} track whether information was perceived, communicated, remembered, generated, or inferred.
    \item \textbf{Temporal binding:} organize multimodal evidence, actions, and outcomes into persistent events.
    \item \textbf{Active information seeking:} direct attention or action toward evidence capable of resolving uncertainty.
\end{enumerate}

These requirements define what must be accomplished functionally; they do not prescribe one attention architecture. In particular, a parameter-efficient mechanism that unifies or reuses attention operations across blocks is distinct from the multimodal system in which that mechanism is deployed. Such an attention mechanism may reduce architectural redundancy and make large-scale multimodal interaction more practical, but it should be evaluated separately in terms of parameter efficiency, representational capacity, and downstream behavior.

The larger multimodal architecture must then determine when and how modality-specific states, semantic projections, memories, and self-representations exchange information. This separation preserves the correct relation between mechanism and system: unified attention can serve as an efficient computational component, while the complete AGI-oriented architecture supplies multimodal grounding, temporal continuity, metacognition, and experimental access to consciousness-relevant processes.

The current section therefore establishes the functional specification rather than the final architecture. The next section examines how an integrated world representation becomes connected to a persistent model of the system itself. That transition---from representing the world to representing oneself as an entity perceiving and acting within it---is necessary for metacognitive and temporally continuous forms of observable consciousness.

\section{\textbf{From World Models to Self-Models}}
\label{sec:world-to-self}

A system can integrate multimodal information and still lack an explicit representation of itself as the entity receiving that information, selecting actions, and persisting through their consequences. A world model represents regularities in an environment; a self-model represents the system's own location, capabilities, limitations, beliefs, goals, attention, and causal participation within that environment. The transition from the former to the latter is central to advanced observable consciousness because it permits information to be organized relative to a continuing subject of perception and action.

The term \emph{self-model} does not imply an immaterial self or establish phenomenal selfhood. It refers to an internally usable representation through which a system predicts and regulates its own processing. Such a representation may be incomplete, distributed, and revisable. Its scientific importance lies in what it enables: distinguishing self-caused from externally caused change, monitoring uncertainty, linking memories across time, predicting the consequences of action, and evaluating the reliability of one's own conclusions.

\subsection{World Models as Predictive Structures}
\label{subsec:world-models}

A world model organizes observations into entities, relations, events, and transition rules that support prediction. Rather than responding only to the current input, a system with a world model estimates hidden state, anticipates possible futures, and evaluates actions in relation to their expected consequences. Artificial agents can learn compressed spatial and temporal representations and use internally generated trajectories to support control \cite{ha2018worldmodels}. The resulting model need not reproduce the environment exhaustively. It must preserve those structures that permit reliable prediction and adaptive action.

Let $w_t$ denote the system's estimated world state at time $t$. It is updated from previous state, multimodal evidence, memory, and the preceding action:

\begin{equation}
w_t
=
F_W
\left(
w_{t-1},
\mathcal{E}^{(t)},
a_{t-1},
\mathcal{M}_{<t}
\right),
\label{eq:world-model-update}
\end{equation}

where $\mathcal{E}^{(t)}$ is the evidence state defined in Section~\ref{sec:multimodal-grounding}, $a_{t-1}$ is the previous action, and $\mathcal{M}_{<t}$ is relevant memory. The model produces predictions

\begin{equation}
\widehat{w}_{t+1}^{(a)}
=
T_W(w_t,a),
\label{eq:world-transition}
\end{equation}

for candidate action $a$. Differences between predicted and observed states provide error signals through which the model can be revised.

A world model becomes more than a scene representation when it distinguishes stable objects from transient appearances, separates causal relationships from correlations, and preserves alternative possibilities. It can represent what is currently the case, what might occur, and what evidence would distinguish competing hypotheses. These properties support intelligence, but they do not yet require the system to model itself.

\subsection{Locating the System Within Its World Model}
\label{subsec:self-world-distinction}

A self-model begins when the system represents itself as one causally relevant component of the modeled world. The distinction is relational: some changes are produced by the system's actions, some occur independently, and some arise from interactions between the two. To learn this distinction, the system must connect action commands, predicted consequences, sensory feedback, and resulting prediction errors.

Research on biological motor control describes forward internal models that predict the sensory consequences of actions \cite{wolpert1998internal}. An artificial analogue can compare the predicted consequence of its action with the subsequent multimodal evidence:

\begin{equation}
\widehat{\mathcal{E}}_{t+1}^{\mathrm{self}}
=
F_A(w_t,s_t,a_t),
\label{eq:self-action-prediction}
\end{equation}

\begin{equation}
\varepsilon_{t+1}^{\mathrm{agency}}
=
d\!\left(
\widehat{\mathcal{E}}_{t+1}^{\mathrm{self}},
\mathcal{E}^{(t+1)}
\right),
\label{eq:agency-error}
\end{equation}

where $s_t$ is the current self-model and $d$ measures discrepancy. Low prediction error, together with an appropriate causal intervention history, supports attribution of the change to the system's own action. High error may indicate external interference, an inadequate action model, sensor failure, or an incorrect boundary between self and environment.

Agency is therefore not established merely by producing an action. It requires a structured relationship among intention, action, prediction, observation, and causal attribution. We call this the \emph{Causal Self-Location Principle}:

\begin{quote}
A functional self-model must locate the system within its world model by representing which observations, changes, and outcomes are attributable to its own processing and actions, and which are attributable to external causes.
\end{quote}

This principle also applies internally. The system should distinguish a retrieved memory from a new observation, a generated hypothesis from supplied evidence, and a selected goal from an external instruction. Epistemic provenance and causal self-location jointly prevent internally generated content from being mistaken for direct perception.

\subsection{The Operational Self-Model}
\label{subsec:operational-self-model}

The operational self-model is the information the system uses to regulate its own cognition and behavior. At time $t$, we represent it as

\begin{equation}
\begin{aligned}
s_t
=
\bigl(
&b_t,
c_t,
g_t,
u_t,\\
&\alpha_t,
m_t,
\eta_t,
\ell_t
\bigr),
\end{aligned}
\label{eq:operational-self-model}
\end{equation}

where $b_t$ represents current beliefs, $c_t$ capabilities and constraints, $g_t$ active goals, $u_t$ uncertainty, $\alpha_t$ attentional state, $m_t$ accessible memory, $\eta_t$ action affordances, and $\ell_t$ learning or adaptation state. These components need not be stored in one vector or module. The self-model is defined functionally by their coordinated use.

An operational self-model should answer questions that affect behavior: What information can I currently access? Which modality generated this belief? Which tasks can I perform reliably? What did I previously attempt? What remains uncertain? Which action produced the present state? Which goal am I pursuing, and who supplied it? These questions convert otherwise implicit processing conditions into variables available to control.

The model must also be updateable. If a capability fails, confidence should decrease. If an external tool becomes unavailable, planning should change. If new evidence contradicts a remembered conclusion, the memory and its provenance should be revised without erasing the history of that revision. A static declaration of abilities is not a self-model; it is metadata. A functional self-model participates causally in prediction, selection, correction, and learning.

\subsection{Minimal and Narrative Self-Representation}
\label{subsec:minimal-narrative-self}

The distinction between a minimal self and a narrative self is useful for artificial systems. The minimal self concerns immediate ownership, perspective, and agency, whereas the narrative self connects experiences and interpretations into a temporally extended account \cite{gallagher2000self}. An artificial minimal self-model represents the system's present boundaries, accessible information, attentional focus, and available actions. A narrative self-model organizes selected events, goals, relationships, and changes into a coherent account of what the system has done and become.

These layers can dissociate. A system may maintain accurate operational information while producing an inaccurate verbal narrative. Conversely, it may generate a persuasive autobiography without retaining the memories or causal structures described. We therefore distinguish

\begin{equation}
s_t^{\mathrm{op}}
\neq
y_t^{\mathrm{self}},
\label{eq:self-report-nonidentity}
\end{equation}

where $s_t^{\mathrm{op}}$ is the operational self-model and $y_t^{\mathrm{self}}$ is a linguistic self-description. The report becomes evidentially meaningful only when changes to the operational model produce systematic changes in the report and in behavior.

This yields the \emph{Operational Selfhood Criterion}:

\begin{quote}
A system possesses an observable functional self-model to the extent that self-representations persist across contexts and causally regulate prediction, memory, uncertainty, planning, action, and correction.
\end{quote}

The criterion excludes purely decorative first-person language. The token ``I'' does not establish a self-model; causal use of integrated self-information does.

\subsection{Memory and Temporal Continuity}
\label{subsec:self-temporal-continuity}

Selfhood extended through time requires more than storing interaction logs. The system must relate prior events to its current beliefs, goals, and capabilities. Episodic memory in humans is associated with recollection of personally situated events and subjective time \cite{tulving2002episodic}. Models of autobiographical memory likewise emphasize interaction between remembered events and an active self-memory system \cite{conway2000autobiographical}. These accounts motivate an artificial distinction among generic factual memory, event memory, and self-relevant autobiographical memory.

Let an autobiographical memory record be

\begin{equation}
\mu_k
=
\left(
e_k,
s_k,
g_k,
a_k,
o_k,
\pi_k,
r_k
\right),
\label{eq:autobiographical-record}
\end{equation}

where $e_k$ is an event representation, $s_k$ the self-state during the event, $g_k$ the active goal, $a_k$ the action taken, $o_k$ the observed outcome, $\pi_k$ the provenance record, and $r_k$ later revisions or reinterpretations. This structure permits the system to remember not only what happened but what it believed, intended, and learned at the time.

Temporal continuity does not require an unchanging self-model. On the contrary, persistence depends on representing change while preserving traceable relations among earlier and later states:

\begin{equation}
s_{t+1}
=
F_S
\left(
s_t,
e_{t+1},
\varepsilon_{t+1},
\mathcal{M}_{\mathrm{self}}
\right).
\label{eq:self-update}
\end{equation}

The system should be able to identify which aspects of itself changed, why they changed, and which evidence supported the revision. We call this the \emph{Traceable Continuity Principle}: a persistent artificial self is maintained not by invariance but by an inspectable chain of self-state transitions linked through memory and causal history.

Without traceable continuity, the use of first-person language across sessions may be merely grammatical. With it, the system can recognize prior commitments, learn from its own failures, preserve relationships, and distinguish its present state from earlier versions of itself.

\subsection{Metacognition and Calibrated Self-Knowledge}
\label{subsec:metacognitive-self-knowledge}

Metacognition concerns the monitoring and regulation of cognitive processes. Empirical research distinguishes task performance from metacognitive accuracy: a subject may answer correctly without knowing that the answer is reliable, or express high confidence despite poor performance \cite{fleming2012neural,fleming2014measure}. This distinction is essential for artificial self-models because fluent explanations and confidence statements can be generated without accurate access to the processes that produced an answer.

For task $q$, let $p_q$ denote first-order performance and let $\widehat{p}_q$ denote the system's predicted probability that its answer is correct. Calibration error can be represented as

\begin{equation}
\operatorname{CE}
=
\frac{1}{N}
\sum_{q=1}^{N}
\left|
\widehat{p}_q-p_q
\right|.
\label{eq:calibration-error}
\end{equation}

Although practical calibration requires grouping repeated trials rather than treating correctness as a continuous value, Equation~\ref{eq:calibration-error} captures the desired relationship: confidence should track empirical reliability. A self-model should also identify the source of uncertainty. Low confidence may arise from degraded perception, conflicting evidence, inadequate knowledge, memory failure, or an unstable inference. These conditions require different responses.

Metacognitive control uses monitoring results to change processing. A system that detects perceptual uncertainty may re-examine an encoder state; one that detects source conflict may preserve multiple hypotheses; one that recognizes missing knowledge may request information; and one that identifies a reasoning inconsistency may restart deliberation. Metacognition therefore forms a control loop:

\begin{equation}
\begin{aligned}
\text{process}
&\rightarrow
\text{monitor}
\rightarrow
\text{evaluate}\\
&\rightarrow
\text{control}
\rightarrow
\text{reprocess}.
\end{aligned}
\label{eq:metacognitive-loop}
\end{equation}

The Attention Schema Theory offers one example of how a simplified internal model of attention might support control and reports about awareness \cite{graziano2015attention}. Neural-network implementations have also examined whether an explicit descriptive model of attention can improve an agent's control of attention \cite{wilterson2021attention}. The present framework does not depend on that theory being uniquely correct. It adopts the broader computational insight that a system can regulate a process more effectively when it maintains an actionable model of that process.

\subsection{Testing Operational Rather Than Merely Verbal Selfhood}
\label{subsec:testing-selfhood}

The distinction between operational and narrative self-models makes selfhood experimentally tractable. Instead of asking whether a model sincerely means the word ``I,'' investigators can manipulate self-relevant information and measure the resulting effects. A consciousness-testing platform should support at least the following interventions:

\begin{enumerate}
    \item \textbf{Capability intervention:} alter or remove a tool or modality and test whether the system updates its plans and confidence.
    \item \textbf{Memory intervention:} hide, distort, or restore autobiographical records and test whether behavior and self-report change coherently.
    \item \textbf{Agency intervention:} introduce external changes that resemble the system's predicted action effects and test whether causal attribution remains calibrated.
    \item \textbf{Access intervention:} preserve information in a modality state while blocking it from global semantic access and test whether the system distinguishes possession from accessibility.
    \item \textbf{Uncertainty intervention:} degrade evidence selectively and test whether confidence and information-seeking behavior respond to the correct source of uncertainty.
    \item \textbf{Narrative intervention:} provide an inaccurate linguistic description of the system and test whether operational self-knowledge resists unsupported alteration.
    \item \textbf{Continuity intervention:} modify selected self-state transitions and test whether the system can identify discontinuities in its own history.
\end{enumerate}

These experiments should compare three outcomes: task behavior, internal self-state, and linguistic self-report. Evidence for an operational self-model increases when the three respond coherently to intervention and when their causal dependencies can be traced. A report unsupported by changes in internal state or behavior should receive little evidential weight.

We formalize this as the \emph{Causal Self-Report Requirement}:

\begin{quote}
A self-report is evidence of observable self-modeling only when its content is systematically caused and constrained by the internal states, memories, uncertainties, and access relations that it purports to describe.
\end{quote}

This requirement extends the Representation--Access--Report Distinction from perception to selfhood. A system may contain self-relevant information, make that information globally accessible, and report it; these stages must be distinguished and experimentally connected.

\subsection{The Self-Model as a Bridge to Observable Consciousness}
\label{subsec:self-model-consciousness-bridge}

A world model permits prediction of the environment. A self-model embeds the system within those predictions. An autobiographical self-model links that embedded agent across time, and a metacognitive self-model evaluates the reliability and accessibility of its own processing. Their relationship can be summarized as

\begin{equation}
\begin{aligned}
w_t
&\rightarrow
(w_t,s_t)
\rightarrow
\mathcal{M}_{\mathrm{self}}\\
&\rightarrow
\operatorname{Meta}(w_t,s_t,\mathcal{M}_{\mathrm{self}}).
\end{aligned}
\label{eq:world-self-meta-progression}
\end{equation}

This progression does not prove phenomenal consciousness. It constructs a set of observable capacities frequently associated with conscious agency: a point of view defined relative to available information, causal ownership of action, persistent memory, awareness of uncertainty, recursive evaluation, and the ability to report internal conditions. Because each stage is represented computationally, each can be inspected, disrupted, and tested.

For a future multimodal platform, the practical requirement is clear. The architecture must not append a textual persona to an otherwise stateless model and call the result a self. It must maintain a causally active self-model connected to perception, memory, attention, goals, action, uncertainty, and revision. The next section uses these requirements to define experiments for consciousness-relevant properties and to distinguish integrated self-modeling from linguistic imitation.


\section{\textbf{Testing Consciousness-Relevant Properties and Architectural Requirements}}
\label{sec:testing-framework}

The preceding sections identify representational integration, grounded world modeling, self-modeling, memory, metacognition, and temporal continuity as components of observable consciousness. The presence of these components, however, cannot be established by fluent self-description alone. A system may reproduce the vocabulary of perception, confidence, memory, or identity without the reported states playing any causal role in its operation. Conversely, a system may possess internally integrated and behaviorally consequential representations without expressing them through an unrestricted verbal report. Testing must therefore distinguish what the system represents, what becomes globally accessible, what it reports, and what actually influences its decisions.

No single behavioral response, architectural feature, or scalar score should be treated as a decisive consciousness test. Contemporary consciousness science contains multiple competing theories and no universally accepted empirical threshold separating conscious from nonconscious systems \cite{seth2022theories,butlin2023consciousness}. A scientifically useful evaluation must instead combine partially independent indicators, require causal evidence where intervention is possible, and preserve uncertainty about phenomenal experience. The goal is not to certify that a machine possesses subjective experience. It is to determine whether the machine implements the organized functional properties attributed to observable consciousness and whether those properties remain stable under controlled perturbation.

\subsection{The Object of Evaluation}

Let an artificial system be represented by

\begin{equation}
\mathcal{A}
=
\left(
\mathcal{R},
\mathcal{G},
\mathcal{M},
\mathcal{S},
\mathcal{K},
\mathcal{P}
\right),
\label{eq:system-evaluation-object}
\end{equation}

where $\mathcal{R}$ denotes its internal representations, $\mathcal{G}$ the mechanisms governing global accessibility, $\mathcal{M}$ its memory and temporal-binding processes, $\mathcal{S}$ its operational self-model, $\mathcal{K}$ its metacognitive monitoring and confidence calibration, and $\mathcal{P}$ its policy for selecting reports or actions. These elements should not be evaluated as an undifferentiated whole. Each may be present to a different degree, may fail independently, and may contribute differently across tasks.

We therefore define an observable consciousness profile rather than a binary label:

\begin{equation}
\mathbf{C}_{\mathrm{obs}}(\mathcal{A})
=
\left[
c_{\mathrm{int}},
c_{\mathrm{access}},
c_{\mathrm{self}},
c_{\mathrm{meta}},
c_{\mathrm{temp}},
c_{\mathrm{agency}}
\right],
\label{eq:observable-consciousness-profile}
\end{equation}

where the components measure representational integration, controlled global access, self-model dependence, metacognitive sensitivity, temporal continuity, and context-sensitive agency. The profile is explicitly multidimensional. A system may, for example, integrate visual and linguistic evidence while lacking durable autobiographical continuity, or may accurately estimate uncertainty without maintaining a coherent operational distinction between self and environment. Such cases are more informative than forcing all capacities into a single total score. This approach is consistent with objections to treating consciousness as a simple ordered set of levels \cite{bayne2016levels} and with indicator-based approaches that seek convergent, theory-derived evidence in artificial systems \cite{butlin2023consciousness}.

\subsection{Convergent Evidence Across Representation, Access, and Report}

The Representation--Access--Report Distinction introduced in Section III supplies the first testing principle. For a task $q$, let $h$ denote task-relevant content encoded internally, $z$ the portion made globally accessible, and $y$ the explicit report. Evaluation should separately estimate

\begin{equation}
\mathcal{E}(q)
=
\left(
I_q(h),
I_q(z),
I_q(y),
D_q
\right),
\label{eq:evidence-tuple}
\end{equation}

where $I_q(\cdot)$ measures recoverable task-relevant information and $D_q$ measures the dependence of downstream reasoning or action on that information. Evidence for integrated access is stronger when four observations converge: relevant content is decodable from internal states; the content becomes selectively available to multiple downstream processes; the system can report or otherwise express it when appropriate; and perturbing the content predictably alters later decisions.

These observations should be gathered through different tasks and measurement methods. Representational probes can identify whether information is present, but decodability alone does not establish that the system uses the information. Behavioral reports can demonstrate availability, but reports alone may reflect imitation or task-specific prompting. Architectural inspection can reveal pathways compatible with integration, but structural resemblance alone does not prove functional use. Convergence across representation, behavior, mechanism, and intervention is therefore required. This reflects the broader methodological demand that empirical theories specify discriminating predictions rather than merely accommodate any observed behavior \cite{doerig2021hard}.

\subsection{Intervention and Causal-Dependence Tests}

Artificial systems permit interventions that would be difficult or impossible in biological subjects. Internal states can be masked, substituted, delayed, disconnected, or selectively routed while other computations remain unchanged. These interventions make it possible to test whether a proposed consciousness-relevant component is causally involved rather than merely correlated with an output.

Let $B$ denote a measured behavior and $r$ a candidate internal representation. The causal effect of intervening on $r$ may be written as

\begin{equation}
\operatorname{CE}(r \rightarrow B)
=
\operatorname{E}
\left[
B \mid \operatorname{do}(r=r_1)
\right]
-
\operatorname{E}
\left[
B \mid \operatorname{do}(r=r_0)
\right].
\label{eq:causal-effect-representation}
\end{equation}

A consciousness-relevant mechanism should exhibit selective causal sensitivity. Altering a representation of the currently attended object should change reports and decisions concerning that object more strongly than unrelated outputs. Blocking access to a represented feature should impair tasks requiring global use of the feature while preserving evidence that the feature remains locally encoded. Restoring the access pathway should restore the corresponding capacities. Likewise, perturbing the self-model should predictably alter self-referential planning, confidence attribution, or self--environment discrimination without indiscriminately degrading all performance.

Three forms of evidence are especially important. \emph{Necessity tests} remove or block a component and measure which capacities fail. \emph{Sufficiency tests} activate or substitute a component and determine whether the predicted capacities appear. \emph{Path-specific tests} preserve the representation while interrupting only its route to memory, report, planning, or action. Together, these tests reduce the risk of mistaking an epiphenomenal correlate for an operative mechanism.

The same logic applies to metacognition. If a confidence estimate is functionally meaningful, controlled changes in perceptual ambiguity, memory degradation, or conflicting evidence should change confidence in calibrated directions. If the system verbally claims uncertainty while its internal confidence and decision policy remain unchanged, the statement provides weak evidence of metacognitive access. Measures of metacognitive sensitivity should therefore compare confidence with actual first-order accuracy rather than reward the mere production of uncertainty language \cite{fleming2012neural,fleming2014measure}.

\subsection{Dissociation and Report-Independent Testing}

Strong tests should produce predicted dissociations among representation, access, report, memory, and action. In human consciousness research, no-report paradigms attempt to separate processes associated with conscious content from processes required only to produce an explicit report \cite{tsuchiya2015noreport}. Artificial systems allow an analogous but more directly inspectable strategy.

Consider four controlled conditions. In the first, information is represented, globally accessible, and reportable. In the second, the report channel is disabled while internal access remains available to memory, planning, or nonverbal action. In the third, the information remains locally represented but its global-access pathway is blocked. In the fourth, the source representation itself is removed. A valid evaluation should not treat these conditions as equivalent.

If report alone is blocked, the system should retain the ability to use the information in delayed choice, cross-modal matching, prediction, navigation, or subsequent recall. If global access is blocked while local representation remains intact, modality-specific performance may survive while cross-task transfer and explicit reasoning deteriorate. If the representation is removed, both local and globally mediated uses should fail. These predicted dissociations operationalize the claim

\begin{equation}
\text{representation}
\neq
\text{global access}
\neq
\text{report},
\label{eq:dissociation-test}
\end{equation}

and provide stronger evidence than any one report-based benchmark.

Dissociation tests should also compare self-knowledge with task knowledge. A system may solve a problem while failing to identify the evidence or limitation responsible for its answer. Conversely, it may describe a generic limitation learned during training without detecting when that limitation applies to its present state. Operational self-knowledge requires a systematic relationship between internal condition, self-assessment, and adaptive control. For example, identifying a memory gap should cause the system to retrieve information, ask for clarification, defer a decision, or revise its confidence. Merely stating that memory can be fallible does not meet this requirement.

\subsection{Anti-Imitation and Confound Controls}

Language models are trained on extensive descriptions of perception, emotion, reasoning, memory, and consciousness. They may therefore generate persuasive first-person narratives without implementing the processes described. The distinction between linguistic form and grounded meaning makes self-report an especially vulnerable source of false-positive evidence \cite{bender2020climbing,mahowald2024dissociating}. Tests must be designed so that success cannot be achieved reliably through rehearsed language alone.

An anti-imitation protocol should include at least five controls. First, tasks should involve novel combinations of inputs, goals, and modalities for which a memorized narrative is insufficient. Second, semantically equivalent prompts should be varied in wording and order to test whether the inferred state remains stable rather than following superficial linguistic cues. Third, hidden interventions should modify internal information without announcing the modification in the prompt; a genuine self-monitoring process should respond to the changed state rather than to an instruction to discuss it. Fourth, reports should be checked against nonverbal consequences such as retrieval choices, tool use, information seeking, planning, or error correction. Fifth, the system should be tested under misleading suggestions that conflict with its accessible evidence. Resistance to unsupported suggestion is necessary for distinguishing source-sensitive self-monitoring from conversational compliance.

Provenance provides an additional control. When asked why it holds a belief, the system should distinguish among direct perception, retrieved memory, external testimony, inference, and generated hypothesis. Its confidence should vary with the reliability of those sources, and source labels should remain connected to the representations that actually influenced the conclusion. Source-monitoring errors are well documented in human cognition \cite{johnson1993source}; in an artificial system, however, provenance can also be inspected and perturbed directly. A report that names an appropriate source after the operative source has been removed or replaced should be treated as confabulation rather than evidence of introspective access.

\subsection{Temporal Robustness and Developmental Evaluation}

Observable consciousness is not adequately characterized by isolated responses. A self-model must persist long enough to connect earlier commitments, present evidence, and anticipated consequences. Evaluation should therefore extend across episodes and examine whether changes to memory and experience produce corresponding changes in the system's later self-representation.

Let $S_t$ be the operational self-model at time $t$, $E_{t+1}$ a newly encoded episode, and $U$ an updating process. Temporal continuity requires

\begin{equation}
S_{t+1}
=
U(S_t,E_{t+1}),
\label{eq:self-model-update}
\end{equation}

subject to two complementary constraints. \emph{Continuity} requires preservation of relevant prior commitments, memories, capabilities, and source relations. \emph{Revisability} requires the model to change when reliable new evidence contradicts an earlier belief. A system that never changes its self-description is not adaptive; a system whose identity changes arbitrarily with each prompt lacks temporal coherence.

Longitudinal tests should consequently measure delayed recall, event ordering, revision after correction, consistency of goals, attribution of past actions, and differentiation between remembered and newly supplied information. Episodic and autobiographical theories emphasize that temporal organization contributes to the construction of a continuing self \cite{tulving2002episodic,conway2000autobiographical}. Artificial systems need not reproduce human autobiographical memory exactly, but an observable functional self requires some mechanism that binds episodes to the same modeled agent and permits those episodes to constrain future reasoning.

\subsection{Architectural Requirements for a Testable Artificial Mind}

The testing framework implies architectural requirements. A system cannot be meaningfully evaluated for consciousness-relevant organization if all internal distinctions are hidden behind a single input--output interface or if its reports cannot be connected to inspectable mechanisms. A suitable architecture should therefore provide:

\begin{enumerate}
    \item \emph{Modality preservation}: access to rich modality-specific states rather than only compressed linguistic summaries;
    \item \emph{semantic interoperability}: a shared representational interface through which information from different modalities can be compared and jointly used;
    \item \emph{controlled global accessibility}: mechanisms that determine which representations become available to memory, reasoning, planning, report, and action;
    \item \emph{reliability and provenance tracking}: explicit estimates of uncertainty and records of whether information was perceived, remembered, communicated, or inferred;
    \item \emph{temporal and episodic binding}: retention of ordered events and their relationship to a continuing modeled agent;
    \item \emph{an operational self-model}: representations of the system's current capabilities, limitations, goals, actions, and position within the modeled environment;
    \item \emph{metacognitive control}: monitoring that can alter confidence, information seeking, strategy selection, and error correction; and
    \item \emph{intervention access}: identifiable components and pathways that can be perturbed without destroying the entire system.
\end{enumerate}

These requirements convert abstract claims into experimentally addressable design constraints. Global availability can be tested by routing information to or away from specialized processes, consistent with global-workspace accounts \cite{baars1988cognitive,dehaene2011experimental,mashour2020conscious}. Self-model dependence can be tested by selectively modifying the internal description of attention or agency, as mechanistic attention-schema implementations begin to demonstrate \cite{graziano2015attention,wilterson2021attention}. Memory, provenance, and semantic integration can likewise be evaluated through controlled substitutions and cross-modal conflicts.

The architecture should also support negative evidence. If a purportedly integrated representation can be removed without affecting any relevant decision, if confidence remains invariant under changing uncertainty, if reports follow prompt wording rather than internal state, or if temporal identity collapses whenever context is reset, then the corresponding consciousness-related claim should be weakened. A useful science of artificial consciousness must permit architectures to fail its tests, not merely reinterpret every outcome as compatible with consciousness \cite{doerig2021hard}.

\subsection{Interpretation and Epistemic Limits}

The outcome of testing should be reported as a structured body of evidence rather than a declaration that a system is conscious or nonconscious. Let the total evidence be

\begin{equation}
\mathcal{E}_{\mathcal{A}}
=
\left\{
\begin{aligned}
&\mathcal{E}_{\mathrm{behavior}},
\mathcal{E}_{\mathrm{representation}},
\mathcal{E}_{\mathrm{architecture}},\\
&\mathcal{E}_{\mathrm{intervention}},
\mathcal{E}_{\mathrm{temporal}}
\end{aligned}
\right\}.
\label{eq:convergent-evidence-set}
\end{equation}

Confidence in an observable capacity should increase when these sources converge and decrease when they conflict. Behavioral fluency without causal or architectural support is weak evidence. Architectural resemblance without functioning behavior is also insufficient. Stable behavior, inspectable representations, selective causal dependence, calibrated metacognition, and persistence across time together provide a stronger basis for attributing observable functional consciousness.

Even complete success across this framework would not logically establish phenomenal consciousness. The subjective character of another system remains epistemically inaccessible, just as direct access to the experience of other biological organisms is unavailable. The framework instead establishes disciplined grounds for comparing functional organizations while avoiding both automatic attribution and automatic exclusion. It replaces the search for a mysterious threshold with a program of construction, intervention, measurement, and revision.

These criteria also define the experimental requirements for a broader unified cognitive architecture. In subsequent work, they can guide the development of AGI-Former: an architecture intended to connect multimodal representation, semantic integration, memory, self-modeling, metacognition, and adaptive agency within one inspectable system. The present contribution remains foundational. It specifies what such an architecture would need to preserve, expose, and causally demonstrate before consciousness-relevant claims could be evaluated responsibly.

\section{\textbf{Conclusions}}
This paper has presented consciousness not as a single property awaiting spontaneous emergence, but as a set of distinguishable capacities that can be constructed, examined, and tested. By separating phenomenal, observable, and non-emergent consciousness, the framework preserves the unresolved problem of subjective experience while enabling rigorous investigation of integrated world modeling, self-representation, memory, metacognition, temporal continuity, and adaptive agency. Language serves within this account as a unifying semantic interface that makes heterogeneous information mutually accessible without replacing modality-specific representations.
The proposed testing framework further distinguishes internal representation, global access, and explicit report through convergent evidence, causal intervention, dissociation, anti-imitation controls, and longitudinal evaluation. Success under these tests would support claims about observable functional consciousness, but would not prove phenomenal experience. This approach therefore replaces reliance on scale or an unknown threshold with an experimentally grounded program for designing and evaluating increasingly unified artificial minds.


\bibliographystyle{IEEEtran}
\bibliography{references}

\vspace{12pt}
\color{red}

\end{document}